% ============================================================
%  Chapitre 1 — État de l'art
% ============================================================

\chapter{État de l'art}
\label{chap:etat_art}

\section{Notions fondamentales}

\subsection{Intelligence Artificielle}

L'Intelligence Artificielle (IA) désigne l'ensemble des techniques permettant
à une machine de simuler des capacités cognitives humaines telles que la
perception, le raisonnement et l'apprentissage. Dans le domaine du traitement
de l'information, l'IA regroupe des sous-disciplines allant de la logique
symbolique aux systèmes d'apprentissage automatique.

\subsection{Machine Learning}

Le Machine Learning (ML) est une branche de l'IA dans laquelle un système
apprend à partir de données, sans être explicitement programmé pour chaque
tâche. Plutôt que de définir des règles manuellement, un algorithme de ML
ajuste ses paramètres internes en minimisant une fonction de coût sur un
ensemble de données étiquetées.

Les principales familles d'algorithmes de ML comprennent :
\begin{itemize}
  \item les méthodes basées sur les arbres de décision (Random Forest,
        Gradient Boosting) ;
  \item les machines à vecteurs de support (SVM) ;
  \item les réseaux de neurones artificiels.
\end{itemize}

\subsection{Deep Learning}

Le Deep Learning (DL) est un sous-domaine du ML qui utilise des réseaux de
neurones profonds — c'est-à-dire composés de nombreuses couches — pour
apprendre des représentations hiérarchiques à partir des données. Chaque
couche transforme la représentation de la couche précédente pour extraire des
caractéristiques de plus en plus abstraites.

Le DL a démontré des performances remarquables sur des tâches de vision par
ordinateur, de traitement du langage naturel et de traitement de séries
temporelles, devenant la méthode de référence pour de nombreuses applications.

\subsection{Réseau de Neurones Convolutionnel (CNN)}

Un Réseau de Neurones Convolutionnel (CNN) est une architecture de DL
spécialisée dans le traitement des données structurées spatialement (images,
séries temporelles). Il repose sur l'opération de convolution, qui applique
un filtre glissant sur l'entrée pour détecter des motifs locaux.

Dans le cadre des séries temporelles à une dimension, une convolution 1D
de noyau de taille $k$ sur une séquence de longueur $T$ et de $C$ canaux
produit une nouvelle séquence où chaque position intègre l'information de
$k$ pas de temps voisins :

\begin{equation}
  y[t] = \sum_{i=0}^{k-1} w[i] \cdot x[t+i]
  \label{eq:conv1d}
\end{equation}

Les CNN capturent efficacement les motifs locaux — pic de NDVI, dates de
semis — mais leur champ réceptif est limité à la taille du noyau.

\subsection{Transformer}

Le Transformer est une architecture introduite par Vaswani et al. \cite{vaswani2017}
fondée sur le mécanisme d'attention. Contrairement aux CNN, le Transformer
capture directement les dépendances à longue portée entre tous les éléments
d'une séquence, indépendamment de leur distance.

Le cœur du Transformer est le mécanisme de \textbf{Multi-Head Self-Attention}
(MHSA), qui calcule pour chaque position un vecteur de contexte pondéré par
l'attention portée à toutes les autres positions :

\begin{equation}
  \text{Attention}(Q, K, V) = \text{softmax}\!\left(\frac{QK^\top}{\sqrt{d_k}}\right) V
  \label{eq:attention}
\end{equation}

où $Q$, $K$, $V$ sont les matrices de requêtes, clés et valeurs, et $d_k$
est la dimension des clés.

L'encodage positionnel est nécessaire pour injecter l'ordre de la séquence,
le MHSA étant invariant à la permutation.

\subsection{Imagerie satellitaire Sentinel-2}

Sentinel-2 est une constellation de deux satellites (2A et 2B) développée
par l'ESA dans le cadre du programme Copernicus. Elle acquiert des images
multispectrales avec 13 bandes, dont 10 à une résolution spatiale de 10 m
ou 20 m. La période de revisite est de 5 jours à l'équateur.

Pour ce projet, 10 bandes spectrales sont utilisées sur une saison agricole
échantillonnée en 36 pas de temps, soit un tenseur de forme $(N, 10, 36)$
par région d'étude.

\section{Méthodes de classification de cultures}

La classification de cultures par télédétection a évolué des méthodes
traditionnelles vers les approches de DL au cours de la dernière décennie.

\subsection{Méthodes classiques}

Les méthodes classiques — Random Forest \cite{breiman2001}, SVM — traitent
chaque pixel indépendamment à partir de statistiques de la série temporelle
(moyenne, variance, indices végétatifs). Elles sont interprétables mais
requièrent une ingénierie des caractéristiques manuelle et échouent à
modéliser les dynamiques temporelles fines.

\subsection{Approches récurrentes (RNN/LSTM)}

Les réseaux récurrents (LSTM, GRU) ont été les premiers à exploiter
directement la dimension temporelle des séries Sentinel-2 \cite{russwurm2018}.
Cependant, ils souffrent de la disparition du gradient sur de longues
séquences et de leur complexité en entraînement.

\subsection{Approches CNN temporels}

Des architectures CNN purement temporelles (TempCNN \cite{pelletier2019})
modélisent la phénologie des cultures via des convolutions 1D. Elles sont
plus légères que les RNN mais leur champ réceptif reste local.

\subsection{Approches Transformer}

Le succès des Transformers en NLP a motivé leur adaptation aux séries
temporelles satellitaires. PSE+TAE \cite{garnot2021} et LTAE introduisent
des mécanismes d'attention sur la dimension temporelle. Ces modèles
capturent bien les tendances globales mais sont coûteux en paramètres.

\subsection{Approches hybrides CNN-Transformer}

Des architectures hybrides combinent la capacité des CNN à extraire des
motifs locaux et celle des Transformers à modéliser les dépendances globales.
MCTNet \cite{wang2024} appartient à cette catégorie : ses blocs CTFusion
traitent la série temporelle en parallèle par un CNN et un Transformer, puis
fusionnent leurs sorties par concaténation.

\section{Tableau comparatif des approches}

\begin{table}[H]
\centering
\caption{Comparaison des principales approches de classification de cultures}
\label{tab:comparaison_approches}
\begin{tabular}{lcccc}
\toprule
\textbf{Méthode} & \textbf{Temporel} & \textbf{Global} & \textbf{Léger} & \textbf{OA (Arkan.)} \\
\midrule
Random Forest     & \textcolor{red}{Non}   & \textcolor{red}{Non}   & \textcolor{green!60!black}{Oui} & $\sim$87\%  \\
LSTM              & \textcolor{green!60!black}{Oui}  & \textcolor{green!60!black}{Oui}  & \textcolor{red}{Non}   & $\sim$90\%  \\
TempCNN           & \textcolor{green!60!black}{Oui}  & \textcolor{red}{Non}   & \textcolor{green!60!black}{Oui} & $\sim$92\%  \\
PSE+TAE           & \textcolor{green!60!black}{Oui}  & \textcolor{green!60!black}{Oui}  & \textcolor{red}{Non}   & $\sim$93\%  \\
\textbf{MCTNet}   & \textcolor{green!60!black}{Oui}  & \textcolor{green!60!black}{Oui}  & \textcolor{green!60!black}{Oui} & \textbf{95\%+} \\
\bottomrule
\end{tabular}
\end{table}

MCTNet se distingue par son compromis entre capacité de modélisation et
efficacité paramétrique, avec environ 56 000 paramètres seulement.

\cleardoublepage
